% Seções finais do relatório (implementação, calibração, resultados e
% conclusão). Pensado para ser incorporado ao relatório parcial; compila
% também isoladamente. Tabelas geradas por scripts/make_tables.py.
\documentclass[11pt,a4paper]{article}
\usepackage[T1]{fontenc}
\usepackage[utf8]{inputenc}
\usepackage[brazil]{babel}
\usepackage{lmodern}
\usepackage[margin=2.5cm]{geometry}
\usepackage{amsmath,amssymb}
\usepackage{booktabs}
\usepackage{graphicx}
\usepackage{pgfplots}
\pgfplotsset{compat=1.18}
\usepackage[round]{natbib}
\usepackage[hidelinks]{hyperref}

\newcommand{\topdois}{\operatorname{top}_2}

\begin{document}

\section*{Alterações em relação à proposta}

A implementação segue a proposta, com três diferenças. Primeiro, a geração de
vizinhos não enumera todos os movimentos viáveis: um movimento é sorteado por
rejeição, o que preserva exatamente a distribuição uniforme entre os
movimentos viáveis (Seção~\ref{sec:impl}). Segundo, o modelo resolvido pelo
solver substitui a família (PLI-3) por uma linearização compacta
equivalente (Seção~\ref{sec:solver}). Terceiro, a calibração da
Seção~\ref{sec:calib} alterou as probabilidades-alvo $p_0$ e $p_{\min}$ e as
probabilidades dos operadores; $\alpha$, $L$ e $K_{\max}$ foram mantidos.

\section{Implementação}\label{sec:impl}

O Simulated Annealing \citep{Kirkpatrick1983} foi implementado em C++17 sem bibliotecas externas e
compilado com g++~15.2 (\texttt{-O2 -march=native}). Os experimentos foram
executados em um Intel Core i5-10210U (4 núcleos físicos, 8 \emph{threads},
até 4{,}2\,GHz) com 7\,GB de memória, sob Ubuntu~26.04. Para reduzir o tempo
total, até oito processos foram executados simultaneamente (três instâncias do
solver e cinco do SA); os tempos relatados são tempos de relógio de cada
processo sob essa carga e, portanto, conservadores.

\paragraph{Estruturas de dados.} Além do vetor \texttt{in\_solution}, de
$E(C)$ e de $A(C)$, a solução mantém o conjunto $C$ e a fronteira
$N(C)=\{v\notin C : v \text{ tem vizinho em } C\}$ em vetores com índice
reverso, que permitem inserir, remover e sortear um elemento em $O(1)$, e o
número de vizinhos em $C$ de cada setor. Uma inclusão ou remoção de $u$ custa
$O(\deg u)$. Em vez dos dois maiores valores de $f$, são mantidos os
\emph{três} maiores: com eles, $W$ após uma remoção ou troca é obtido em
$O(1)$, e a recomputação em $O(|C|)$ só ocorre quando um desses três setores
deixa a solução em um movimento aceito.

\paragraph{Conectividade.} Os candidatos do \emph{Add} são os setores da
fronteira, logo $C\cup\{v\}$ é sempre conexo. Uma remoção de $u\neq p$
preserva a conectividade se, e somente se, $u$ não é ponto de articulação de
$G[C]$. Os pontos de articulação são obtidos pela busca em profundidade de
\citet{Tarjan1972} a partir do Palácio, em $O(|C|+\sum_{u\in C}\deg u)$, e
recalculados sob demanda após cada movimento aceito. Para uma troca $(u,v)$,
se $u$ não é articulação basta que $v$ tenha vizinho em $C\setminus\{u\}$.
Se $u$ é articulação, as componentes de $G[C]-u$ são a que contém o Palácio
e uma para cada filho $c$ de $u$ na árvore de busca com
$\mathrm{low}(c)\geq\mathrm{disc}(u)$; $(C\setminus\{u\})\cup\{v\}$ é conexo
se, e somente se, $v$ tem vizinho em todas elas. A componente de cada
vizinho de $v$ é identificada pelo intervalo de tempos de descoberta da
subárvore que o contém, com uma busca binária entre os filhos de $u$. O teste
custa $O(\deg v\cdot\log\deg u)$ e dispensa a BFS da proposta. A rotina foi
validada contra uma BFS em cerca de $1{,}2$ milhão de casos aleatórios.

\paragraph{Sorteio dos movimentos.} A enumeração completa dos movimentos
viáveis (Seção~5 da proposta) custaria $O(|C|\cdot|N(C)|)$ testes por troca.
Em seu lugar, sorteia-se uniformemente um candidato em um superconjunto dos
movimentos viáveis (a fronteira para o \emph{Add}, $C$ para o \emph{Remove} e
$C\times N(C)$ para o \emph{Swap}) e rejeita-se o candidato inviável, até
64, 64 e 4096 tentativas, respectivamente; esgotadas as tentativas, os
candidatos são enumerados. Como cada tentativa é uniforme e independente, o
primeiro candidato aceito é uniforme entre os viáveis, e a enumeração de
reserva também o é; a distribuição dos vizinhos é, portanto, exatamente a da
proposta. Ao final, a melhor solução é verificada por um procedimento
independente (BFS e recálculo de $W$), e um segundo verificador, em Python,
confere todos os arquivos de solução.

\section{Modelo no solver}\label{sec:solver}

O modelo da Seção~1.4 foi resolvido com o HiGHS~1.15.1
\citep{HuangfuHall2018}, com uma \emph{thread}, \texttt{mip\_rel\_gap}$=0$ e
limite de 1800\,s por instância. Escrita explicitamente, a família (PLI-3)
tem $1+n+\binom{n}{2}$ restrições com até $n+2$ coeficientes cada: para
$n=2000$, cerca de $2\cdot10^6$ linhas e $4\cdot10^9$ coeficientes, o que
excede a memória disponível. Ela foi substituída pela linearização usual da
soma dos $k$ maiores valores \citep{OgryczakTamir2003}, com variáveis
auxiliares $t\geq0$ e $s_i\geq0$:
\begin{align}
  \sum_{i\in S} a_i x_i + 2t + \sum_{i\in S} s_i &\leq H, \label{eq:cap}\\
  s_i &\geq f_i x_i - t, \qquad \forall i\in S. \label{eq:si}
\end{align}
Para $x$ fixo, $\min_{t\geq0}\, 2t+\sum_i \max\{0, f_i x_i - t\}$ é
exatamente a soma dos dois maiores valores $f_i x_i$, atingida com $t$ igual
ao segundo maior valor. Logo existe $(t,s)$ que satisfaz
\eqref{eq:cap}--\eqref{eq:si} se, e somente se, $x$ satisfaz (PLI-3), e o
conjunto de soluções viáveis em $x$ não muda. O modelo resultante tem
$2n+2|E|+1$ variáveis e $2n+4|E|+1$ restrições; para jifc12, são 33\,027
variáveis, 62\,053 restrições e 188\,155 coeficientes.

\section{Calibração dos parâmetros}\label{sec:calib}

Os parâmetros foram calibrados um fator por vez nas instâncias jifc01, jifc05
e jifc09 ($n=500$, uma de cada família), com sementes 101 a 103 (101 a 108
na última etapa), distintas das usadas nos experimentos finais. A medida de
qualidade é o desvio percentual médio para o BKV,
$100(\mathrm{BKV}-\mathrm{SF})/\mathrm{BKV}$, negativo quando a solução
supera o BKV. Como na proposta, as temperaturas são definidas pela
probabilidade de aceitar uma piora típica \citep{Johnson1989}. Como $T_0$ e
$T_{\min}$ são ambos proporcionais à escala $d$, o
número de níveis independe da instância: é o menor $k$ com
$\alpha^k\leq \ln p_0/\ln p_{\min}$. Na proposta ($p_0=0{,}8$,
$p_{\min}=10^{-3}$, $\alpha=0{,}95$) são sempre 67 níveis, de modo que
$K_{\max}=100$ nunca é atingido.

\begin{table}[htbp]
  \centering
  \caption{Desvio médio para o BKV (\%) em função de $p_0$ e $p_{\min}$, com
    $\alpha=0{,}95$, $L=50|S|$ e operadores 40/30/30. Entre parênteses, o
    número de níveis de temperatura.}
  \label{tab:temp}
  \begin{tabular}{lrrrr}
\toprule
& \multicolumn{4}{c}{$p_{\min}$} \\ \cmidrule(l){2-5}
$p_0$ & $10^{-3}$ & $10^{-10}$ & $10^{-20}$ & $10^{-40}$ \\ \midrule
0{,}8 & 2{,}08 (67) & 0{,}46 (91) & 0{,}25 (104) & 0{,}24 (118) \\
0{,}5 & 1{,}94 (45) & 0{,}09 (69) & $-$0{,}11 (82) & $-$0{,}12 (96) \\
0{,}3 & 1{,}80 (35) & 0{,}18 (58) & 0{,}04 (72) & 0{,}02 (85) \\
\bottomrule
\end{tabular}

\end{table}

A Tabela~\ref{tab:temp} mostra que a temperatura final é o fator dominante.
Com $p_{\min}=10^{-3}$, o algoritmo termina com taxa de aceitação de 4\% a
17\% e a solução corrente ainda abaixo da solução gulosa; em jifc09, nenhuma
das três execuções melhorou a solução inicial. Reduzir $p_{\min}$ para
$10^{-20}$ diminui o desvio médio entre 1{,}8 e 2{,}1 pontos percentuais;
reduções adicionais trazem ganho desprezível. A temperatura inicial com
$p_0=0{,}8$ destrói a solução gulosa sem benefício: $p_0=0{,}5$ obteve
resultados melhores com menos níveis. Foram escolhidos $p_0=0{,}5$ e
$p_{\min}=10^{-20}$, o que resulta em 82 níveis, abaixo de $K_{\max}$.

\begin{table}[htbp]
  \centering
  \caption{Desvio médio para o BKV (\%) em função de $\alpha$ e $L$, com
    $p_0=0{,}5$, $p_{\min}=10^{-20}$ e operadores 40/30/30.}
  \label{tab:resf}
  \begin{tabular}{lrrrr}
\toprule
& & \multicolumn{3}{c}{$L/|S|$} \\ \cmidrule(l){3-5}
$\alpha$ & níveis & 25 & 50 & 100 \\ \midrule
0{,}90 & 40 & 0{,}92 & 0{,}71 & 0{,}06 \\
0{,}95 & 82 & 0{,}73 & $-$0{,}11 & $-$0{,}30 \\
0{,}98 & 208 & 0{,}01 & $-$0{,}22 & $-$0{,}39 \\
\bottomrule
\end{tabular}

\end{table}

O esforço de uma execução é proporcional ao número de níveis vezes $L$. Na
Tabela~\ref{tab:resf}, entre as configurações de esforço semelhante
($\alpha=0{,}90$ com $L=100|S|$, $\alpha=0{,}95$ com $L=50|S|$ e
$\alpha=0{,}98$ com $L=25|S|$), a da proposta é a melhor.
Aumentar o esforço em até cinco vezes reduz o desvio em no máximo 0{,}3 ponto
percentual, e
por isso $\alpha=0{,}95$ e $L=50|S|$ foram mantidos.

\begin{table}[htbp]
  \centering
  \caption{Desvio médio para o BKV (\%) por mistura de operadores, com 8
    sementes, e fração das iterações sem vizinho viável.}
  \label{tab:ops}
  \begin{tabular}{lrrrrr}
\toprule
Add/Remove/Swap (\%) & jifc01 & jifc05 & jifc09 & média & sem vizinho (\%) \\ \midrule
40/30/30 & 0{,}49 & $-$1{,}74 & 1{,}73 & 0{,}16 & 35{,}1 \\
25/25/50 & 0{,}40 & $-$2{,}54 & 1{,}31 & $-$0{,}28 & 21{,}3 \\
20/20/60 & 0{,}37 & $-$2{,}49 & 1{,}10 & $-$0{,}34 & 17{,}1 \\
15/25/60 & 0{,}22 & $-$2{,}99 & 0{,}83 & $-$0{,}65 & 11{,}9 \\
10/10/80 & 0{,}31 & $-$2{,}71 & 1{,}21 & $-$0{,}40 & 8{,}7 \\
10/20/70 & 0{,}17 & $-$3{,}00 & 0{,}60 & $-$0{,}75 & 7{,}8 \\
\bottomrule
\end{tabular}

\end{table}

A Tabela~\ref{tab:ops} explica o efeito das probabilidades dos operadores.
As boas soluções saturam a capacidade do aqueduto, e nelas quase nenhuma
inclusão é viável: com 40\% de \emph{Add}, 35\% das iterações terminam sem
vizinho e são desperdiçadas. Deslocar probabilidade para o \emph{Swap}, que
substitui setores sem alterar muito o consumo, e para o \emph{Remove}, que
libera capacidade, reduz esse desperdício para menos de 8\% e melhorou o
resultado nas três instâncias. A configuração final usa \emph{Add},
\emph{Remove} e \emph{Swap} com probabilidades 10\%, 20\% e 70\%.

\section{Resultados}\label{sec:res}

A Tabela~\ref{tab:res} apresenta, para cada instância, a média de dez
replicações (sementes 1 a 10) do SA com a configuração calibrada e o resultado
do solver. SI é o valor da solução gulosa, SF o melhor valor encontrado pelo
SA, ``mel.'' a melhoria percentual média de SF sobre SI,
$100(\mathrm{SF}-\mathrm{SI})/\mathrm{SI}$, e ``desv.'' o desvio percentual
médio para o BKV. Valores de SF máximo acima do BKV estão em negrito.

\begin{table}[htbp]
  \centering
  \footnotesize
  \setlength{\tabcolsep}{3.5pt}
  \caption{Resultados do SA (configuração calibrada) e do HiGHS.}
  \label{tab:res}
  \begin{tabular}{lrrrrrrrrrrrr}
\toprule
& \multicolumn{8}{c}{Simulated Annealing (10 replicações)} & \multicolumn{4}{c}{HiGHS (1800\,s)} \\ \cmidrule(lr){2-9} \cmidrule(l){10-13}
Inst. & SI & SF méd. & dp & SF máx. & BKV & mel.\,(\%) & desv.\,(\%) & t\,(s) & valor & limite & gap\,(\%) & t\,(s) \\ \midrule
jifc01 & 82\,135 & 87\,854 & 142 & 88\,009 & 88\,063 & 6{,}96 & 0{,}24 & 15{,}4 & 88\,063 & 88\,063 & 0{,}00 & 1777 \\
jifc02 & 129\,877 & 141\,871 & 284 & \textbf{142\,211} & 136\,257 & 9{,}23 & $-$4{,}12 & 52{,}2 & 406 & 146\,310 & 99{,}72 & 1800 \\
jifc03 & 169\,270 & 195\,850 & 287 & \textbf{196\,232} & 186\,418 & 15{,}70 & $-$5{,}06 & 132{,}0 & 700 & 203\,302 & 99{,}66 & 1800 \\
jifc04 & 201\,738 & 236\,294 & 336 & \textbf{236\,933} & 209\,722 & 17{,}13 & $-$12{,}67 & 176{,}8 & 868 & 245\,385 & 99{,}65 & 1802 \\
jifc05 & 56\,018 & 64\,206 & 115 & \textbf{64\,381} & 62\,273 & 14{,}62 & $-$3{,}10 & 10{,}0 & 63\,284 & 66\,890 & 5{,}39 & 1800 \\
jifc06 & 95\,883 & 105\,606 & 114 & \textbf{105\,805} & 103\,748 & 10{,}14 & $-$1{,}79 & 35{,}2 & 542 & 109\,035 & 99{,}50 & 1801 \\
jifc07 & 108\,947 & 125\,639 & 378 & \textbf{126\,222} & 118\,878 & 15{,}32 & $-$5{,}69 & 67{,}2 & 106 & 132\,456 & 99{,}92 & 1810 \\
jifc08 & 132\,435 & 155\,151 & 239 & \textbf{155\,492} & 141\,621 & 17{,}15 & $-$9{,}55 & 131{,}3 & 131 & 162\,454 & 99{,}92 & 1808 \\
jifc09 & 44\,338 & 45\,069 & 85 & 45\,202 & 45\,345 & 1{,}65 & 0{,}61 & 15{,}2 & 45\,256 & 45\,706 & 0{,}98 & 1800 \\
jifc10 & 57\,410 & 58\,979 & 176 & 59\,284 & 59\,368 & 2{,}73 & 0{,}65 & 54{,}9 & 461 & 60\,233 & 99{,}23 & 1801 \\
jifc11 & 74\,736 & 77\,983 & 185 & \textbf{78\,343} & 76\,498 & 4{,}34 & $-$1{,}94 & 108{,}0 & 364 & 80\,900 & 99{,}55 & 1800 \\
jifc12 & 86\,477 & 90\,699 & 181 & \textbf{91\,073} & 88\,180 & 4{,}88 & $-$2{,}86 & 176{,}5 & 377 & 93\,497 & 99{,}60 & 1809 \\
\midrule
Média & & & & & & 9{,}99 & $-$3{,}77 & & & & & \\
\bottomrule
\end{tabular}

\end{table}

Todas as 120 soluções produzidas pelo SA foram confirmadas pelos dois
verificadores. O SA melhorou a solução gulosa em todas as execuções, em média
10{,}0\%, com ganhos maiores na família jifc01--jifc04 (7{,}0\% a 17{,}1\%) e
menores na família jifc09--jifc12 (1{,}7\% a 4{,}9\%), em que a solução
gulosa já é boa. O desvio padrão entre replicações não passa de 0{,}3\% do
valor médio, de modo que o resultado depende pouco da semente.

Em nove das doze instâncias, a média do SA supera o BKV, e a vantagem cresce
com o tamanho da instância dentro de cada família, chegando a 12{,}7\% em
jifc04. Nas três instâncias restantes (jifc01, jifc09 e jifc10), a melhor
replicação fica entre 0{,}06\% e 0{,}32\% abaixo do BKV, e a média, entre
0{,}24\% e 0{,}65\%. Em jifc01, o solver provou que o BKV é ótimo, e a melhor
solução do SA está a 0{,}06\% dele. Como os BKV das instâncias maiores foram
superados por margem grande, os limites superiores do solver são uma
referência mais informativa: nenhuma solução do SA os ultrapassa, e a melhor
solução de cada instância está entre 0{,}06\% e 4{,}7\% abaixo do limite.
Esse valor limita a distância ao ótimo por cima, pois inclui também a folga
do próprio limite.

O desempenho do solver, com gap calculado como
$100(\text{limite}-\text{valor})/\text{limite}$, depende fortemente do
tamanho da instância. Com $n=500$, ele resolveu jifc01 até a otimalidade em
1777\,s e terminou jifc05 e jifc09 com gaps de 5{,}4\% e 1{,}0\%. Em jifc05,
sua solução supera o BKV, mas fica 1{,}7\% abaixo da melhor do SA; em jifc09,
é a única instância em que o solver encontrou solução melhor que todas as
replicações do SA (45\,256 contra 45\,202), ainda abaixo do BKV. Com
$n\ge1000$, as soluções inteiras obtidas em 1800\,s contêm apenas o Palácio e
poucos vizinhos, com esplendor entre 106 e 868, e o gap fica acima de 99\%;
em jifc04, jifc07, jifc08, jifc11 e jifc12, a árvore de busca não passou de
dois nós. A causa é a relaxação linear fraca do modelo de fluxo: com
$M=n-1$ em (PLI-6)--(PLI-7), um setor com $x_i$ próximo de zero pode
receber e repassar fluxo, de modo que a relaxação quase ignora a
conectividade e o limite se aproxima do de uma mochila fracionária. Um modelo
mais forte exigiria, por exemplo, restrições de corte de conectividade
geradas sob demanda \citep{MagnantiWolsey1995}, o que está fora do escopo
deste trabalho.

Esse comportamento do solver é coerente com o padrão dos BKV: ótimo ou
próximo do ótimo nas instâncias menores e cada vez mais distante dele nas
maiores, o que sugere que os BKV destas últimas foram obtidos com esforço
limitado.

O tempo do SA cresce com $n$ e com $|C|$, de 10\,s nas instâncias com
$n=500$ a 177\,s nas com $n=2000$, e é sempre muito menor que o tempo dado ao
solver. Em média, a melhor solução é encontrada após 82\% a 99\% do tempo de
execução, isto é, nos níveis finais e frios, como ilustra a
Figura~\ref{fig:conv}. No primeiro nível, em que cerca de 70\% dos movimentos
são aceitos, a solução corrente cai bem abaixo da gulosa; a partir daí, ela
sobe de forma contínua, supera a solução inicial por volta do nível 30 e se
estabiliza perto do nível 70, quando menos de 0{,}1\% dos movimentos ainda é
aceito.

A Tabela~\ref{tab:prop} compara as duas configurações nas mesmas sementes.
A configuração da proposta termina ainda quente e fica, em média, 3{,}6 pontos
percentuais atrás da calibrada; em toda a família jifc09--jifc12, nenhuma de
suas replicações melhorou a solução gulosa. A configuração calibrada é melhor
em todas as instâncias, e não apenas nas três usadas para calibrar, e é também
mais rápida, apesar de ter mais níveis: com 40\% de \emph{Add}, muitas
iterações esgotam as tentativas de rejeição e recorrem à enumeração da
fronteira.

\begin{table}[htbp]
  \centering
  \small
  \caption{Configuração da proposta e configuração calibrada (sementes 1 a 5).}
  \label{tab:prop}
  \begin{tabular}{lrrrrrr}
\toprule
& \multicolumn{3}{c}{Proposta} & \multicolumn{3}{c}{Calibrada} \\ \cmidrule(lr){2-4} \cmidrule(l){5-7}
Inst. & SF méd. & desv.\,(\%) & t\,(s) & SF méd. & desv.\,(\%) & t\,(s) \\ \midrule
jifc01 & 86\,668 & 1{,}58 & 24{,}6 & 87\,826 & 0{,}27 & 15{,}3 \\
jifc02 & 139\,449 & $-$2{,}34 & 78{,}9 & 141\,803 & $-$4{,}07 & 52{,}5 \\
jifc03 & 191\,844 & $-$2{,}91 & 167{,}9 & 195\,869 & $-$5{,}07 & 140{,}1 \\
jifc04 & 231\,920 & $-$10{,}58 & 292{,}4 & 236\,367 & $-$12{,}70 & 175{,}7 \\
jifc05 & 61\,066 & 1{,}94 & 15{,}8 & 64\,167 & $-$3{,}04 & 11{,}0 \\
jifc06 & 100\,663 & 2{,}97 & 94{,}6 & 105\,676 & $-$1{,}86 & 33{,}6 \\
jifc07 & 118\,383 & 0{,}42 & 176{,}3 & 125\,531 & $-$5{,}60 & 64{,}4 \\
jifc08 & 146\,119 & $-$3{,}18 & 343{,}3 & 155\,058 & $-$9{,}49 & 130{,}9 \\
jifc09 & 44\,338 & 2{,}22 & 33{,}9 & 45\,088 & 0{,}57 & 16{,}6 \\
jifc10 & 57\,410 & 3{,}30 & 96{,}2 & 58\,969 & 0{,}67 & 54{,}1 \\
jifc11 & 74\,736 & 2{,}30 & 205{,}1 & 78\,006 & $-$1{,}97 & 109{,}3 \\
jifc12 & 86\,477 & 1{,}93 & 320{,}6 & 90\,635 & $-$2{,}78 & 180{,}0 \\
\midrule
Média & & $-$0{,}20 & & & $-$3{,}76 & \\
\bottomrule
\end{tabular}

\end{table}

\begin{figure}[htbp]
  \centering
  \begin{tikzpicture}
    \begin{axis}[width=0.85\textwidth, height=6cm, xlabel={nível de temperatura $k$},
                 ylabel={esplendor}, legend pos=south east, xmin=0, xmax=81,
                 scaled y ticks=false,
                 y tick label style={/pgf/number format/.cd, fixed, 1000 sep={\,}}]
      \addplot[thin, gray] table[x=level, y=current, col sep=comma]
        {../../results/sa_calibrada/trace/jifc06_calibrada_s1.csv};
      \addplot[thick, black] table[x=level, y=best, col sep=comma]
        {../../results/sa_calibrada/trace/jifc06_calibrada_s1.csv};
      \legend{solução corrente, melhor solução}
    \end{axis}
  \end{tikzpicture}
  \caption{Evolução do SA em jifc06 (semente 1), ao final de cada nível.}
  \label{fig:conv}
\end{figure}

\section{Conclusão}

O Simulated Annealing proposto, com amostragem uniforme dos movimentos
viáveis por rejeição e testes de conectividade baseados em pontos de
articulação, produziu soluções viáveis de alta qualidade para todas as
instâncias, superando o BKV em nove das doze e ficando a menos de 0{,}7\% dele
nas demais, em no máximo três minutos por execução. O modelo PLIM, com a
linearização compacta da restrição hídrica, é competitivo nas instâncias com
500 setores, em que provou a otimalidade de jifc01 e superou o SA em jifc09,
mas não produziu soluções úteis em 30 minutos a partir de 1000 setores. Ainda
assim, seus limites superiores situam as soluções do SA a no máximo 4{,}7\%
do ótimo.

A calibração mostrou que a escolha da temperatura final é decisiva neste
problema: a configuração da proposta encerrava a busca antes que a solução
convergisse. Mostrou também que a capacidade hídrica saturada torna as
inclusões quase sempre inviáveis, o que favorece um peso maior para as trocas.
Como a melhor solução ainda é encontrada nos últimos níveis, uma linha de
continuação é estender a fase fria, por exemplo com reaquecimentos ou uma
busca local final. Do lado exato, cortes de conectividade gerados sob demanda
devem fortalecer a relaxação e permitir avaliar melhor a distância ao ótimo.

\begin{thebibliography}{9}
\bibitem[Huangfu e Hall(2018)]{HuangfuHall2018}
Q.~Huangfu e J.~A.~J. Hall. Parallelizing the dual revised simplex method.
\emph{Mathematical Programming Computation}, 10(1):119--142, 2018.

\bibitem[Johnson et~al.(1989)]{Johnson1989}
D.~S. Johnson, C.~R. Aragon, L.~A. McGeoch e C.~Schevon. Optimization by
simulated annealing: an experimental evaluation; part I, graph partitioning.
\emph{Operations Research}, 37(6):865--892, 1989.

\bibitem[Kirkpatrick et~al.(1983)]{Kirkpatrick1983}
S.~Kirkpatrick, C.~D. Gelatt e M.~P. Vecchi. Optimization by simulated
annealing. \emph{Science}, 220(4598):671--680, 1983.

\bibitem[Magnanti e Wolsey(1995)]{MagnantiWolsey1995}
T.~L. Magnanti e L.~A. Wolsey. Optimal trees. Em M.~O. Ball et~al.
(eds.), \emph{Network Models}, Handbooks in Operations Research and
Management Science, vol.~7, p. 503--615. Elsevier, 1995.

\bibitem[Ogryczak e Tamir(2003)]{OgryczakTamir2003}
W.~Ogryczak e A.~Tamir. Minimizing the sum of the $k$ largest functions in
linear time. \emph{Information Processing Letters}, 85(3):117--122, 2003.

\bibitem[Tarjan(1972)]{Tarjan1972}
R.~E. Tarjan. Depth-first search and linear graph algorithms. \emph{SIAM
Journal on Computing}, 1(2):146--160, 1972.
\end{thebibliography}

\end{document}
